\documentclass[../../dis.tex]{subfiles}
\begin{document}
\chapter{SQL-1 - SQL Introduction}
\vfill
\section{Relational terminology}
Although we already defined the relational model, let's align our vocabulary with SQL:
\begin{dashlist}
    \item \textbf{Relation} translates to a \textbf{table}.
    \item \textbf{Tuple} translates to a \textbf{row}.
    \item \textbf{Attribute} translates to a \textbf{column}.
\end{dashlist}

Throughout this chapter, we use the following example \texttt{Students} table:

\begin{dbtable}{Students}{lllcc}
\textbf{\underline{sid}} & \textbf{name} & \textbf{login} & \textbf{age} & \textbf{gpa} \\ \hline
53666 & Jones & jones@cs & 18 & 3.4 \\ \hline
53688 & Smith & smith@ee & 18 & 3.2 \\ \hline
53777 & White & white@cs & 19 & 4.0 \\
\end{dbtable}

\begin{concept}{Primary Key (PK)}
The column (or group of columns) that uniquely identifies a row. Two rows \textbf{cannot} share the same PK. Above, \texttt{\underline{sid}} is the primary key.
\end{concept}

\section{Basic SQL queries}

\subsection{The simplest query}
To retrieve all contents of a table, use \texttt{SELECT *}:
\begin{sql}
    SELECT * FROM Students S
\end{sql}
\texttt{S} is an optional \textit{alias} (range variable) for the \texttt{Students} table, useful when referencing columns. This returns the full table shown above.

\subsection{Showing specific columns (Projection)}
To retrieve only certain attributes, replace \texttt{*} with column names:
\begin{sql}
    SELECT S.name, S.login FROM Students S
\end{sql}

\begin{concept}{Projection}
Selecting specific columns from a table. The query above projects \texttt{name} and \texttt{login} from \texttt{Students}.
\end{concept}

\begin{dbtable}{}{ll}
\textbf{name} & \textbf{login} \\ \hline
Jones & jones@cs \\ \hline
Smith & smith@ee \\ \hline
White & white@cs \\
\end{dbtable}

\subsection{Showing specific rows (Selection)}
To filter rows, use the \texttt{WHERE} clause:
\begin{sql}
    SELECT * FROM Students S WHERE S.age = 18
\end{sql}

\begin{concept}{Selection}
Filtering rows based on a condition. The query above selects students whose age equals~18.
\end{concept}

\begin{dbtable}{}{lllcc}
\textbf{\underline{sid}} & \textbf{name} & \textbf{login} & \textbf{age} & \textbf{gpa} \\ \hline
53666 & Jones & jones@cs & 18 & 3.4 \\ \hline
53688 & Smith & smith@ee & 18 & 3.2 \\
\end{dbtable}

\subsection{SQL vocabulary}
The general form of a basic SQL query:
\begin{sql}
    SELECT [DISTINCT] target-list FROM relation-list WHERE qualification
\end{sql}
\begin{dashlist}
    \item \textbf{relation-list}: a list of table names, each possibly followed by an alias.
    \item \textbf{qualification}: comparisons combined with \texttt{AND}, \texttt{OR}, \texttt{NOT}. Each comparison: \texttt{Attr \textit{op} Const} or \texttt{Attr1 \textit{op} Attr2}.
    \item \textbf{target-list}: a list of attributes from the tables in \textit{relation-list}.
    \item \textbf{DISTINCT}: optional keyword; eliminates duplicate rows. By default, SQL keeps duplicates (the result is a \textit{multiset}).
\end{dashlist}

\subsection{Evaluation order}
Conceptually, a SQL query is evaluated as follows:
\begin{enumerate}
    \item \textbf{FROM}: compute the cross-product of all tables.
    \item \textbf{WHERE}: discard tuples that fail the condition (\textit{selection}).
    \item \textbf{SELECT}: keep only the requested columns (\textit{projection}).
    \item If \textbf{DISTINCT} is specified, eliminate duplicate rows.
\end{enumerate}

\section{Querying multiple relations}
To join two tables, list them both in \texttt{FROM} and express the join condition in \texttt{WHERE}. Consider the following \texttt{Enrolled} table alongside \texttt{Students}:

\begin{dbtable}{Enrolled}{llc}
\textbf{\underline{sid}} & \textbf{cid} & \textbf{grade} \\ \hline
53831 & Carnatic101 & C \\ \hline
53831 & Reggae203 & B \\ \hline
53650 & Topology112 & A \\ \hline
53666 & History105 & B \\
\end{dbtable}

``Find all students with grade B'':
\begin{sql}
    SELECT S.name, E.cid
    FROM Students S, Enrolled E
    WHERE S.sid = E.sid AND E.grade = 'B'
\end{sql}

\begin{dbtable}{}{ll}
\textbf{S.name} & \textbf{E.cid} \\ \hline
Jones & History105 \\
\end{dbtable}

\subsection{Na\"ive evaluation in steps}
The evaluation of the query above proceeds conceptually as follows:

\textbf{Step 1 -- Cross Product (FROM):}
Combine every row of \texttt{Students} with every row of \texttt{Enrolled}:

\begin{center}
\scriptsize
\begin{tabular}{|l|l|l|l|c|l|l|c|}
\hline
\rowcolor{DBHeaderBg}
\textbf{S.sid} & \textbf{S.name} & \textbf{S.login} & \textbf{S.age} & \textbf{S.gpa} & \textbf{E.sid} & \textbf{E.cid} & \textbf{E.grade} \\ \hline
53666 & Jones & jones@cs & 18 & 3.4 & 53831 & Carnatic101 & C \\ \hline
53666 & Jones & jones@cs & 18 & 3.4 & 53831 & Reggae203 & B \\ \hline
53666 & Jones & jones@cs & 18 & 3.4 & 53650 & Topology112 & A \\ \hline
53666 & Jones & jones@cs & 18 & 3.4 & 53666 & History105 & B \\ \hline
53688 & Smith & smith@ee & 18 & 3.2 & 53831 & Carnatic101 & C \\ \hline
53688 & Smith & smith@ee & 18 & 3.2 & 53831 & Reggae203 & B \\ \hline
53688 & Smith & smith@ee & 18 & 3.2 & 53650 & Topology112 & A \\ \hline
53688 & Smith & smith@ee & 18 & 3.2 & 53666 & History105 & B \\ \hline
\end{tabular}
\normalsize
\end{center}

\textbf{Step 2 -- Apply the predicate (WHERE):}
Keep only rows where \texttt{S.sid = E.sid} \textbf{and} \texttt{E.grade = 'B'}. Only one row survives: Jones (sid~53666) enrolled in History105 with grade~B.

\textbf{Step 3 -- Project the requested columns (SELECT):}
Keep only \texttt{S.name} and \texttt{E.cid}, yielding the final result: \texttt{(Jones, History105)}.

\section*{Running example: Sailors, Boats, Reserves}
The following three relations are used throughout the rest of this chapter.

\begin{minipage}[t]{0.32\textwidth}
\begin{dbtable}{Sailors}{llcc}
\textbf{\underline{sid}} & \textbf{sname} & \textbf{rating} & \textbf{age} \\ \hline
22 & Dustin & 7 & 45.0 \\ \hline
31 & Lubber & 8 & 55.5 \\ \hline
95 & Bob    & 3 & 63.5 \\
\end{dbtable}
\end{minipage}
\hfill
\begin{minipage}[t]{0.32\textwidth}
\begin{dbtable}{Boats}{llc}
\textbf{\underline{bid}} & \textbf{bname} & \textbf{color} \\ \hline
101 & Interlake & blue  \\ \hline
102 & Interlake & red   \\ \hline
103 & Clipper   & green \\ \hline
104 & Marine    & red   \\
\end{dbtable}
\end{minipage}
\hfill
\begin{minipage}[t]{0.28\textwidth}
\begin{dbtable}{Reserves}{ccl}
\textbf{\underline{sid}} & \textbf{\underline{bid}} & \textbf{day} \\ \hline
22 & 101 & 10/10/96 \\ \hline
95 & 103 & 11/12/96 \\
\end{dbtable}
\end{minipage}

\section{Aggregate operators}

\begin{concept}{Aggregate function}
A function that operates on a set (or multiset) of values from a column and returns a \textbf{single} summary value. Aggregates are a significant extension of relational algebra.
\end{concept}

\begin{center}
\begin{tabular}{@{}ll@{}}
    \toprule
    \textbf{Function} & \textbf{Description} \\
    \midrule
    \texttt{COUNT(*)}              & number of rows \\
    \texttt{COUNT([DISTINCT] A)}   & number of (distinct) non-null values in \texttt{A} \\
    \texttt{SUM([DISTINCT] A)}     & sum of (distinct) values \\
    \texttt{AVG([DISTINCT] A)}     & average of (distinct) values \\
    \texttt{MAX(A)}                & maximum value \\
    \texttt{MIN(A)}                & minimum value \\
    \bottomrule
\end{tabular}
\end{center}

\subsection{Examples}

\textbf{Find the number of sailors:}
\begin{sql}
    SELECT COUNT(*) FROM Sailors S
\end{sql}

\textbf{Find the average age of sailors with rating 10:}
\begin{sql}
    SELECT AVG(S.age) FROM Sailors S WHERE S.rating = 10
\end{sql}

\textbf{Count distinct ratings of sailors named Bob:}
\begin{sql}
    SELECT COUNT(DISTINCT S.rating) FROM Sailors S WHERE S.sname = 'Bob'
\end{sql}

\subsection{Name and age of the oldest sailor(s)}
This example illustrates a common pitfall with aggregates.

\subsubsection*{Incorrect query}
\begin{sql}
    SELECT S.sname, MAX(S.age) FROM Sailors S
\end{sql}
This is \textbf{wrong}: \texttt{MAX} returns a single value over the entire column, so there is no meaningful way to pair it with a particular \texttt{sname}.

\subsubsection*{Correct approach --- subquery}
\begin{sql}
    SELECT S.sname, S.age
    FROM Sailors S
    WHERE S.age = (SELECT MAX(S2.age) FROM Sailors S2)
\end{sql}
The inner query computes the maximum age; the outer query retrieves every sailor matching that value.

\begin{concept}{Scalar subquery}
A subquery that returns exactly one row and one column. It can appear wherever a single value is expected (e.g.\ in a \texttt{WHERE} comparison).
\end{concept}

An equivalent alternative places the subquery on the left:
\begin{sql}
    SELECT S.sname, S.age
    FROM Sailors S
    WHERE (SELECT MAX(S2.age) FROM Sailors S2) = S.age
\end{sql}

\section{GROUP BY}

So far, aggregates have been applied to \emph{all} qualifying tuples at once. Sometimes we need to apply them to each of several \emph{groups} independently.

\textit{Consider: find the age of the youngest sailor for each rating level.}
\begin{dashlist}
    \item We don't know in advance how many rating levels exist or what their values are.
    \item If ratings range from 1 to 10, we \emph{could} write 10 separate queries:
\end{dashlist}

For $i = 1, 2, \ldots, 10$:
\begin{sql}
    SELECT MIN(S.age) FROM Sailors S WHERE S.rating = i
\end{sql}
This is obviously impractical.

\begin{concept}{GROUP BY clause}
Partitions tuples into groups based on equal values of the specified columns, then applies aggregate functions independently to each group.
\end{concept}

\subsection{Using GROUP BY}
We use an extended \texttt{Sailors} table for the following examples:

\begin{dbtable}{Sailors (extended)}{llcc}
\textbf{\underline{sid}} & \textbf{sname} & \textbf{rating} & \textbf{age} \\ \hline
22 & Dustin & 7 & 45.0 \\ \hline
31 & Lubber & 8 & 55.5 \\ \hline
95 & Bob    & 3 & 63.5 \\ \hline
52 & Smith  & 3 & 74.5 \\ \hline
68 & John   & 7 & 56.5 \\ \hline
81 & Ana    & 7 & 71.5 \\
\end{dbtable}

``Find the youngest age for each rating level'':
\begin{sql}
    SELECT S.rating, MIN(S.age) FROM Sailors S GROUP BY S.rating
\end{sql}

The tuples are first partitioned by \texttt{rating} (groups: rating~3, rating~7, rating~8), then \texttt{MIN(S.age)} is evaluated within each group:

\begin{dbtable}{}{cc}
\textbf{rating} & \textbf{MIN(age)} \\ \hline
3 & 63.5 \\ \hline
7 & 45.0 \\ \hline
8 & 55.5 \\
\end{dbtable}

\begin{concept}{GROUP BY semantics}
The \texttt{SELECT} clause may only contain:
\begin{dashlist}
    \item attributes listed in the \texttt{GROUP BY} clause, and
    \item aggregate functions applied to other attributes.
\end{dashlist}
Any non-aggregated column in \texttt{SELECT} that is \emph{not} in \texttt{GROUP BY} is an error.
\end{concept}

\subsection{HAVING --- filtering groups}
\texttt{HAVING} filters \emph{groups} (after grouping), just as \texttt{WHERE} filters \emph{individual tuples} (before grouping).

``Find the youngest age for each rating level that has \textbf{at least two} sailors'':
\begin{sql}
    SELECT S.rating, MIN(S.age)
    FROM Sailors S
    GROUP BY S.rating
    HAVING COUNT(*) > 1
\end{sql}

Groups with only one member (rating~8, containing only Lubber) are eliminated:

\begin{dbtable}{}{cc}
\textbf{rating} & \textbf{MIN(age)} \\ \hline
3 & 63.5 \\ \hline
7 & 45.0 \\
\end{dbtable}

\begin{concept}{HAVING clause}
A filter applied \emph{after} grouping. It uses aggregate conditions to keep or discard entire groups, unlike \texttt{WHERE} which operates on individual rows before any grouping.
\end{concept}

\end{document}
